En general, la complejidad teórica predijo bien el comportamiento
general de estos algoritmos, pero no siempre el rendimiento real.
\texttt{Sort} fue el más parejo entre tiempo y memoria, sin
importar el tipo de entrada; Quick Sort con pivote determinista, en
cambio, mostró justo su punto débil clásico: arreglos ya
ordenados, hasta el punto de no terminar dentro del límite de tiempo,
confirmando de la forma más directa posible su peor caso $O(n^2)$.
Strassen sí rinde más que Naive, pero solo a partir de matrices
razonablemente grandes ($n\geq 256$); para las chicas, su propia
sobrecarga lo deja atrás. Si algo queda claro de este ejercicio es que
la notación asintótica es un buen punto de partida, pero no reemplaza
medir: el contexto de uso y las características concretas de los datos
importan tanto como la complejidad "en el papel".
En consecuencia, la elección de un algoritmo no debería basarse únicamente
en su complejidad asintótica, sino también en las características de los datos, 
el tamaño de la entrada y los costos asociados a su implementación.